% ============================================================================
% MÓDULO 9 - PROYECTO FINAL: TECHRETAIL MÉXICO
% ============================================================================
\chapter{Proyecto final: TechRetail México}
\label{cap:proyecto}

\begin{objetivos}
\begin{itemize}
  \item Plantear un proyecto de Business Intelligence completo a partir de un
        reto de negocio real: preguntas, entregables, cronograma y criterios
        de evaluación.
  \item Construir un proceso ETL auditable: cargar varias fuentes, detectar
        problemas de calidad, documentar cada decisión de limpieza y armar un
        modelo en estrella listo para analizar.
  \item Calcular los KPIs de la empresa y convertir cada tabla o gráfica en
        un hallazgo de negocio con números reales.
  \item Aplicar segmentación de clientes y pronóstico de ventas siendo honesto
        sobre lo que los datos permiten (y lo que no).
  \item Diseñar el tablero ejecutivo, redactar el reporte para la dirección y
        preparar la presentación de resultados.
  \item Organizar el proyecto de forma reproducible y presentarlo en tu
        portafolio para buscar trabajo en BI.
\end{itemize}
\end{objetivos}

Llegaste al final del curso. Ya sabes importar y transformar datos
(Módulo~\ref{cap:manipulacion}), graficarlos con criterio
(Módulo~\ref{cap:visualizacion}), explorarlos con estadística
(Módulo~\ref{cap:eda}), guardarlos en una base de datos
(Módulo~\ref{cap:bases-datos}), construir tableros interactivos
(Módulo~\ref{cap:shiny}), entrenar modelos (Módulo~\ref{cap:ml}) y
automatizar reportes (Módulo~\ref{cap:reportes}). Cada módulo resolvió una
pieza. En la vida real nadie te pide ``una pieza'': te piden \emph{resolver un
problema de negocio}, y tú decides qué herramientas usar, en qué orden y cómo
comunicar lo que encontraste.

Este módulo simula exactamente eso. Vas a trabajar como consultor de BI para
\textbf{TechRetail México}, una cadena de diez sucursales que tiene un año de
datos y muchas preguntas. No vamos a re-enseñar técnicas: cada vez que uses
algo que ya aprendiste, te remitiremos al módulo donde se explicó y nos
concentraremos en lo que distingue a un analista junior de uno senior: el
\emph{razonamiento de negocio}, las \emph{decisiones} (y su justificación) y la
\emph{comunicación} de resultados.

Una advertencia honesta antes de empezar: los datos del curso son simulados.
Eso significa que algunos patrones que esperarías en una tienda real (más
ventas en diciembre, fines de semana más fuertes, clientes premium que gastan
más) pueden no aparecer, o aparecer por azar. Lejos de ser un problema, es un
excelente entrenamiento: aprenderás a decir ``los datos no muestran
diferencia'' con la misma seguridad con la que dirías ``encontramos una
oportunidad de un millón de pesos''. Un buen analista no fuerza conclusiones.

\begin{concepto}{Proyecto integrador de BI}
Un \termino{proyecto integrador} recorre el ciclo completo de la
inteligencia de negocios: \emph{pregunta de negocio} $\rightarrow$
\emph{datos} (ETL) $\rightarrow$ \emph{análisis} (KPIs, exploración,
modelos) $\rightarrow$ \emph{producto} (tablero, reporte) $\rightarrow$
\emph{decisión} (presentación y recomendaciones). El valor no está en
ninguna etapa aislada, sino en que la cadena completa lleve de un dato crudo
a una acción concreta de la dirección.
\end{concepto}

\begin{figure}[H]
\centering
\begin{tikzpicture}[
    fase/.style={draw=primario, fill=primario!6, rounded corners=3pt,
                 minimum width=4.4cm, minimum height=1.45cm, align=center,
                 font=\small\sffamily},
    flecha/.style={-{Stealth[length=2.5mm]}, thick, color=acento}]
  \node[fase] (f1) {\textbf{Fase 1. ETL y modelo}\\ Módulos \ref{cap:manipulacion} y \ref{cap:bases-datos}};
  \node[fase, right=0.9cm of f1] (f2) {\textbf{Fase 2. KPIs y exploración}\\ Módulos \ref{cap:visualizacion} y \ref{cap:eda}};
  \node[fase, right=0.9cm of f2] (f3) {\textbf{Fase 3. Segmentación}\\ \textbf{y pronóstico}\\ Módulo \ref{cap:ml}};
  \node[fase, below=1.1cm of f3] (f4) {\textbf{Fase 4. Dashboard}\\ Módulo \ref{cap:shiny}};
  \node[fase, left=0.9cm of f4] (f5) {\textbf{Fase 5. Reporte ejecutivo}\\ Módulo \ref{cap:reportes}};
  \node[fase, left=0.9cm of f5] (f6) {\textbf{Fase 6. Presentación}\\ Storytelling con datos};
  \draw[flecha] (f1) -- (f2);
  \draw[flecha] (f2) -- (f3);
  \draw[flecha] (f3) -- (f4);
  \draw[flecha] (f4) -- (f5);
  \draw[flecha] (f5) -- (f6);
\end{tikzpicture}
\caption{Las seis fases del proyecto y los módulos del libro que usa cada una.}
\label{fig:m09-fases}
\end{figure}

% ============================================================================
\section{El caso: TechRetail México}
\label{sec:m09-caso}
% ============================================================================

TechRetail México es una cadena minorista con diez sucursales en seis
ciudades del país. Vende productos de diez categorías (desde electrónica y
juguetes hasta alimentos y artículos para mascotas) y registra cada venta en
un sistema de punto de venta. La dirección tiene los datos de 2023 pero nunca
los ha analizado de forma integral: cada gerente de tienda lleva su propio
Excel, y las juntas de resultados terminan en discusiones sobre ``qué número
es el bueno''.

Antes de escribir una sola línea de análisis, prepara el entorno del
proyecto. Fíjate en la \termino{fecha de corte}: es la fecha ``de hoy'' para
el proyecto. El script original del curso usaba \codigo{Sys.Date()}, lo que
hace que los resultados cambien cada día que lo ejecutas (la antigüedad de
una tienda o los días desde la última compra de un cliente serían distintos
mañana). En un proyecto profesional eso es inaceptable: quien revise tu
trabajo dentro de seis meses debe obtener exactamente los mismos números.

\begin{rcode}[title=Configuración del proyecto]
# ---- Paquetes (todos se explicaron en módulos anteriores) ----
library(tidyverse)   # dplyr, ggplot2, readr, tidyr, purrr, lubridate
library(scales)      # formatos de negocio: dollar(), percent(), comma()
library(DBI)         # interfaz a bases de datos (Módulo 5)
library(RSQLite)     # motor SQLite en un archivo local
library(cluster)     # silueta para evaluar k-means (Módulo 7)
library(forecast)    # modelos de pronóstico (Módulo 7)

# ---- Fecha de corte FIJA (nunca Sys.Date() en un proyecto) ----
fecha_corte <- as.Date("2024-01-01")

# ---- Carpetas de salida del proyecto ----
dir.create("resultados", showWarnings = FALSE)
dir.create("datos", showWarnings = FALSE)

# ---- Colores corporativos y tema (los mismos del Módulo 3) ----
color_principal <- "#2a78d6"   # una sola serie
color_resalte   <- "#eb6834"   # lo que queremos destacar
color_gris      <- "#b5b3ad"   # el resto
paleta_libro <- c("#2a78d6", "#eb6834", "#1baf7a", "#eda100",
                  "#e87ba4", "#008300", "#4a3aa7", "#e34948")

tema_libro <- function(base_size = 11) {
  theme_minimal(base_size = base_size) +
    theme(
      plot.title       = element_text(face = "bold", color = "#0b0b0b"),
      plot.subtitle    = element_text(color = "#52514e"),
      plot.caption     = element_text(color = "#52514e", size = rel(0.8)),
      axis.title       = element_text(color = "#52514e"),
      axis.text        = element_text(color = "#52514e"),
      panel.grid.minor = element_blank(),
      panel.grid.major = element_line(color = "#e6e5e1", linewidth = 0.3),
      legend.position  = "top",
      legend.title     = element_text(color = "#52514e")
    )
}

# Formato de pesos sin decimales para textos y tablas
pesos <- function(x) dollar(x, accuracy = 1)
\end{rcode}

Al cargar \paquete{tidyverse} y \paquete{forecast} verás mensajes de
bienvenida o avisos como \codigo{Registered S3 method overwritten}; son
normales y no indican ningún error.

\subsection{La cadena en números}

Describe a la empresa con sus propios datos, no con adjetivos. El archivo
\archivo{tiendas.csv} es el catálogo de sucursales.

\begin{rcode}
tiendas <- read_csv("datasets/tiendas.csv", show_col_types = FALSE)

# Antigüedad de cada sucursal calculada a la fecha de corte
tiendas %>%
  mutate(anios = floor(time_length(
           interval(fecha_apertura, fecha_corte), "years"))) %>%
  select(tienda = nombre_tienda, ciudad, zona, tipo,
         m2 = tamano_m2, empleados, anios) %>%
  arrange(desc(m2))
\end{rcode}

\begin{rsalida}
# A tibble: 10 x 7
   tienda           ciudad      zona   tipo         m2 empleados anios
   <chr>            <chr>       <chr>  <chr>     <dbl>     <dbl> <dbl>
 1 Sucursal Premium CDMX        Sur    Premium     800        40    14
 2 Sucursal Plaza A CDMX        Oeste  Principal   700        35     7
 3 Sucursal Sur     Monterrey   Norte  Estándar    600        30    12
 4 Sucursal Oeste   Puebla      Centro Estándar    550        22    10
 5 Sucursal Centro  CDMX        Centro Principal   500        25    13
 6 Sucursal Plaza B Guadalajara Norte  Estándar    480        21     9
 7 Sucursal Norte   Guadalajara Norte  Estándar    450        20    11
 8 Sucursal Este    CDMX        Este   Estándar    400        18     8
 9 Sucursal Outlet  Tijuana     Norte  Outlet      350        15     5
10 Sucursal Express Querétaro   Centro Express     300        12     3
\end{rsalida}

\begin{rcode}
# Resumen de la red de sucursales
tiendas %>%
  summarise(sucursales = n(),
            ciudades   = n_distinct(ciudad),
            m2_totales = sum(tamano_m2),
            empleados  = sum(empleados),
            m2_por_empleado = round(sum(tamano_m2) / sum(empleados), 1))

tiendas %>% count(ciudad, sort = TRUE)
\end{rcode}

\begin{rsalida}
# A tibble: 1 x 5
  sucursales ciudades m2_totales empleados m2_por_empleado
       <int>    <int>      <dbl>     <dbl>           <dbl>
1         10        6       5130       238            21.6
# A tibble: 6 x 2
  ciudad          n
  <chr>       <int>
1 CDMX            4
2 Guadalajara     2
3 Monterrey       1
4 Puebla          1
5 Querétaro       1
6 Tijuana         1
\end{rsalida}

¿Qué nos dice esto? TechRetail tiene 10 sucursales en 6 ciudades, con
5,130~m\textsuperscript{2} de piso de venta y 238 empleados. La Ciudad de
México concentra 4 sucursales, incluidas la más grande (Sucursal Premium,
800~m\textsuperscript{2} y 40 empleados, la más antigua: 14 años) y la
Sucursal Plaza A. En el otro extremo están la Sucursal Express de Querétaro
(300~m\textsuperscript{2}, abierta en 2020) y la Sucursal Outlet de Tijuana.
Hay cinco formatos (Principal, Estándar, Outlet, Premium y Express), lo cual
será importante: no es justo comparar las ventas absolutas de una tienda de
800~m\textsuperscript{2} con las de una de 300.

\begin{advertencia}
Observa que la ``Sucursal Sur'' está en Monterrey y pertenece a la zona
\emph{Norte}. Los nombres heredados no siempre describen la realidad. En el
análisis usaremos siempre la columna \codigo{zona}, nunca el nombre de la
tienda, para agrupar geográficamente.
\end{advertencia}

\subsection{El reto de la dirección}

\begin{negocio}{El correo de la directora general}
\emph{``Necesitamos dejar de discutir números y empezar a tomar decisiones.
Quiero, en seis semanas: (1) una sola versión de la verdad de nuestras ventas
de 2023, (2) entender qué tiendas, productos y clientes nos hacen ganar
dinero, (3) un tablero que los gerentes consulten cada lunes y (4) una
recomendación clara de en qué enfocarnos en 2024. Preséntalo al comité
directivo en 20 minutos.''}

Traducido a BI: la dirección pide \textbf{calidad y gobierno de datos} (una
sola versión de la verdad), \textbf{análisis descriptivo y diagnóstico}
(qué pasó y por qué), un \textbf{producto de datos recurrente} (tablero) y
\textbf{recomendaciones accionables} (presentación y reporte).
\end{negocio}

\subsection{Las 20 preguntas de negocio}

El primer entregable de cualquier consultoría es la lista de preguntas que
vas a contestar, acordada con el cliente. Evita el ``análisis de todo'': si
una tabla no responde una pregunta, sobra. La
Tabla~\ref{tab:m09-preguntas} agrupa las 20 preguntas en cuatro áreas e
indica dónde se responde cada una en este módulo.

\begin{table}[H]
\centering
\small
\caption{Las 20 preguntas de negocio del proyecto y dónde se responden.}
\label{tab:m09-preguntas}
\begin{tabularx}{\textwidth}{@{}c X l@{}}
\toprule
\# & Pregunta & Dónde se responde \\
\midrule
\multicolumn{3}{@{}l}{\textbf{Ventas}} \\
1 & ¿Cuánto vendimos en 2023, en cuántas transacciones y con qué ticket promedio? & Sec.~\ref{sec:m09-kpis-globales} \\
2 & ¿Cómo evolucionaron las ventas mes a mes? ¿Estamos creciendo? & Sec.~\ref{sec:m09-mensual} \\
3 & ¿Qué categorías generan más ingresos y más margen? & Sec.~\ref{sec:m09-categorias} \\
4 & ¿Qué método de pago prefieren los clientes y cambia el ticket? & Sec.~\ref{sec:m09-pago} \\
5 & ¿Qué ventas podemos esperar en las próximas semanas? & Sec.~\ref{sec:m09-pronostico} \\
\midrule
\multicolumn{3}{@{}l}{\textbf{Clientes}} \\
6 & ¿Cuántos clientes activos tenemos y cuántos no compraron? & Sec.~\ref{sec:m09-kpis-globales} \\
7 & ¿Los clientes premium gastan más que los demás? & Sec.~\ref{sec:m09-premium} \\
8 & ¿Qué segmentos de clientes existen según su comportamiento? & Sec.~\ref{sec:m09-rfm} \\
9 & ¿Cuántos clientes vuelven a comprar y cada cuánto? & Reto~\ref{ej:p-1} \\
10 & ¿Se comportan distinto los clientes según el mes en que llegaron? & Reto~\ref{ej:p-4} \\
\midrule
\multicolumn{3}{@{}l}{\textbf{Productos}} \\
11 & ¿Pocos productos concentran la mayoría de las ventas (Pareto)? & Sec.~\ref{sec:m09-pareto} \\
12 & ¿Qué productos tienen margen negativo en el catálogo? & Sec.~\ref{sec:m09-auditoria} \\
13 & ¿Vendemos al precio de catálogo? & Sec.~\ref{sec:m09-auditoria} \\
14 & ¿Se venden productos marcados como inactivos? & Sec.~\ref{sec:m09-auditoria} \\
15 & ¿Qué productos necesitan reabastecerse pronto? & Reto~\ref{ej:p-3} \\
\midrule
\multicolumn{3}{@{}l}{\textbf{Operaciones}} \\
16 & ¿Qué sucursales y zonas venden más? & Sec.~\ref{sec:m09-tiendas} \\
17 & ¿Qué sucursales son más productivas por m\textsuperscript{2} y por empleado? & Reto~\ref{ej:p-2} \\
18 & ¿Qué días y horas concentran las ventas? & Sec.~\ref{sec:m09-dia-hora} \\
19 & ¿Qué vendedores venden más y se relaciona con su salario? & Sec.~\ref{sec:m09-vendedores} \\
20 & ¿Hay transacciones atípicas que debamos revisar? & Reto~\ref{ej:p-5} \\
\bottomrule
\end{tabularx}
\end{table}

\begin{consejo}
Escribe las preguntas \emph{antes} de abrir los datos y pídele a la dirección
que las priorice. Si al final del proyecto alguien pregunta ``¿y esto para
qué sirve?'', debes poder señalar qué pregunta responde cada gráfica.
\end{consejo}

\subsection{Entregables y cronograma}

El proyecto se organiza en seis fases con un entregable concreto cada una
(Tabla~\ref{tab:m09-cronograma}). Las horas son una estimación para alguien
que termina este libro; en una empresa real, la Fase~1 suele consumir
más tiempo del que todos esperan.

\begin{table}[H]
\centering
\small
\caption{Cronograma del proyecto por fases.}
\label{tab:m09-cronograma}
\begin{tabularx}{\textwidth}{@{}l X X c c@{}}
\toprule
Fase & Actividades principales & Entregable & Semana & Horas \\
\midrule
1. ETL y modelo & Cargar fuentes, auditar calidad, limpiar, modelo en estrella & Datos limpios (RDS y SQLite) y bitácora de decisiones & 1 & 5--6 \\
2. KPIs y exploración & KPIs globales, tendencias, cortes por tienda, producto, cliente y vendedor & Diccionario de KPIs y cuaderno de hallazgos & 2 & 5--6 \\
3. Modelos & Segmentación RFM con k-means y pronóstico de ventas & Segmentos con acciones y pronóstico con intervalos & 3 & 4--5 \\
4. Dashboard & Diseño y construcción del tablero en Shiny & App funcionando con filtros & 4 & 4--5 \\
5. Reporte & Reporte ejecutivo automatizado en R Markdown & PDF o HTML de 1--2 páginas & 5 & 2--3 \\
6. Presentación & Guion, diapositivas y ensayo & Presentación de 8--10 diapositivas & 6 & 2--3 \\
\bottomrule
\end{tabularx}
\end{table}

\subsection{Rúbrica de evaluación}

Si usas este proyecto para un curso, un bootcamp o tu portafolio, evalúalo
con criterios explícitos (Tabla~\ref{tab:m09-rubrica}). Observa que el código
vale menos de la mitad: en BI, un análisis técnicamente perfecto que nadie
entiende no sirve.

\begin{table}[H]
\centering
\small
\caption{Rúbrica de evaluación del proyecto final.}
\label{tab:m09-rubrica}
\begin{tabularx}{\textwidth}{@{}l c X X@{}}
\toprule
Criterio & Peso & Excelente & Insuficiente \\
\midrule
Calidad de datos & 15\% & Auditoría con evidencia, decisiones documentadas y justificadas & Se limpian datos ``a ojo'' o se borran filas sin explicar \\
Código y reproducibilidad & 20\% & Scripts numerados, comentados, fecha de corte fija; corre de principio a fin & Rutas absolutas, \codigo{Sys.Date()}, pasos manuales \\
Análisis y KPIs & 25\% & Cada gráfica responde una pregunta y trae un hallazgo con números & Muchas gráficas sin conclusión \\
Modelos & 10\% & Modelo justificado, validado y con limitaciones explícitas & Modelo sin validar o con conclusiones exageradas \\
Tablero y reporte & 15\% & Tablero claro con KPIs y filtros; reporte breve y automatizado & Tablero saturado; reporte copiado a mano \\
Comunicación & 15\% & Historia clara, recomendaciones accionables con impacto & Lista de datos sin recomendación \\
\bottomrule
\end{tabularx}
\end{table}

\begin{hazlotu}
Agrega tres preguntas propias a la Tabla~\ref{tab:m09-preguntas} (por
ejemplo, sobre marcas o proveedores). Para cada una anota qué archivo y qué
columnas necesitarías para responderla.
\end{hazlotu}

% ============================================================================
\section{Fase 1: ETL y modelo de datos}
\label{sec:m09-etl}
% ============================================================================

La primera fase construye la ``única versión de la verdad'' que pidió la
directora. Seguimos las tres letras del \termino{ETL}: \emph{extraer} (leer
las fuentes), \emph{transformar} (auditar, limpiar, enriquecer) y
\emph{cargar} (guardar el resultado en un formato listo para analizar). Las
funciones de \paquete{dplyr} y \paquete{readr} se explicaron en el
Módulo~\ref{cap:manipulacion} y las bases de datos en el
Módulo~\ref{cap:bases-datos}; aquí nos enfocamos en las decisiones.

\subsection{Extraer: cargar las fuentes}

\begin{rcode}
# Tabla de hechos candidata principal: una fila por transacción
transacciones <- read_csv("datasets/transacciones.csv",
                          show_col_types = FALSE)
# Segunda fuente de ventas (una venta registrada por día)
ventas_retail <- read_csv("datasets/ventas_retail.csv",
                          show_col_types = FALSE)
# Catálogos (futuras dimensiones)
clientes  <- read_csv("datasets/clientes.csv",  show_col_types = FALSE)
productos <- read_csv("datasets/productos.csv", show_col_types = FALSE)
empleados <- read_csv("datasets/empleados.csv", show_col_types = FALSE)
# tiendas ya se cargó en la sección anterior

glimpse(transacciones)
\end{rcode}

\begin{rsalida}
Rows: 1,000
Columns: 14
$ transaccion_id    <dbl> 1, 2, 3, 4, 5, 6, 7, 8, 9, 10, 11, 12, 13,~
$ fecha             <date> 2023-01-13, 2023-01-09, 2023-02-20, 2023-~
$ hora              <time> 13:15:00, 10:45:00, 10:30:00, 10:45:00, 1~
$ cliente_id        <dbl> 155, 155, 75, 51, 173, 135, 90, 168, 13, 9~
$ producto_id       <dbl> 3, 47, 6, 17, 28, 46, 30, 48, 21, 12, 49, ~
$ cantidad          <dbl> 4, 4, 3, 4, 1, 2, 4, 2, 2, 5, 2, 3, 4, 2, ~
$ precio_venta      <dbl> 303.97, 58.07, 275.46, 105.28, 205.47, 534~
$ metodo_pago       <chr> "Tarjeta Crédito", "Tarjeta Débito", "Tarj~
$ tienda_id         <dbl> 10, 2, 7, 6, 7, 7, 8, 8, 5, 2, 7, 6, 8, 2,~
$ vendedor_id       <dbl> 4, 22, 14, 5, 14, 22, 23, 5, 3, 14, 5, 9, ~
$ total_transaccion <dbl> 1215.88, 232.28, 826.38, 421.12, 205.47, 1~
$ mes               <chr> "2023-01", "2023-01", "2023-02", "2023-05"~
$ trimestre         <chr> "Q1", "Q1", "Q1", "Q2", "Q1", "Q1", "Q2", ~
$ dia_semana        <chr> "viernes", "lunes", "lunes", "viernes", "m~
\end{rsalida}

Fíjate en un detalle: en el archivo la hora viene como texto sin cero
inicial (\codigo{"8:00"}, \codigo{"13:15"}), pero \paquete{readr} la
reconoció como hora (\codigo{<time>}). Eso nos permitirá extraer la hora del
día con \codigo{hour()} sin conversiones manuales.

\subsection{¿Qué tabla usar como hechos?}
\label{sec:m09-fuente}

Tenemos \emph{dos} archivos de ventas. Antes de sumar nada, compáralos: si
fueran el mismo registro visto de dos formas, sumarlos duplicaría las ventas;
si son registros distintos, necesitamos saber cuál es el oficial.

\begin{rcode}
# Perfil rápido de cada fuente de ventas
perfil_fuente <- function(datos, nombre, col_total) {
  tibble(fuente   = nombre,
         filas    = nrow(datos),
         columnas = ncol(datos),
         desde    = min(datos$fecha),
         hasta    = max(datos$fecha),
         ventas   = sum(datos[[col_total]]))
}

bind_rows(
  perfil_fuente(transacciones, "transacciones", "total_transaccion"),
  perfil_fuente(ventas_retail, "ventas_retail", "total")
)

# ¿Qué información tiene cada fuente que la otra no?
setdiff(names(transacciones), names(ventas_retail))
setdiff(names(ventas_retail), names(transacciones))

# ¿Coinciden registros (misma fecha, cliente y producto)?
inner_join(transacciones, ventas_retail,
           by = c("fecha", "cliente_id", "producto_id")) %>%
  nrow()
\end{rcode}

\begin{rsalida}
# A tibble: 2 x 6
  fuente        filas columnas desde      hasta       ventas
  <chr>         <int>    <int> <date>     <date>       <dbl>
1 transacciones  1000       14 2023-01-01 2023-12-31 959993.
2 ventas_retail   365       14 2023-01-01 2023-12-31 513917.
[1] "transaccion_id"    "hora"              "precio_venta"     
[4] "metodo_pago"       "total_transaccion"
[1] "precio_unitario" "descuento_pct"   "subtotal"       
[4] "descuento_monto" "total"          
[1] 0
\end{rsalida}

Las dos fuentes cubren todo 2023, pero no son el mismo registro: ninguna
combinación de fecha, cliente y producto coincide entre ellas, y los totales
son muy distintos (\$959,993 contra \$513,917). \textbf{Decisión:} usamos
\archivo{transacciones.csv} como tabla de hechos porque (1) es más granular
(1,000 transacciones contra 365 registros, uno por día), (2) trae
información que la otra no tiene y que la dirección pidió (hora, método de
pago) y (3) su total es consistente con su detalle (lo comprobaremos en la
auditoría). \archivo{ventas\_retail.csv} queda como fuente secundaria: es la
única que registra descuentos, así que la usaremos solo para preguntas sobre
promociones (Reto~\ref{ej:p-6}). \textbf{Nunca sumaremos ambas fuentes.}

\begin{advertencia}
Mezclar fuentes de ventas es el error más caro de un proyecto de BI: el
tablero muestra un total que nadie reconoce y el proyecto pierde credibilidad
en la primera junta. Si dos sistemas registran ventas, documenta cuál es el
oficial para cada pregunta y dilo explícitamente en el reporte.
\end{advertencia}

\subsection{Auditoría de calidad de datos}
\label{sec:m09-auditoria}

Una \termino{auditoría de calidad} revisa sistemáticamente cada tabla antes
de usarla: tamaño, faltantes, duplicados, rangos, integridad de llaves y
reglas de negocio. Conviértela en una función para aplicarla igual a todas
las tablas (y reutilizarla cada mes cuando lleguen datos nuevos).

\begin{rcode}
# Auditoría básica de una tabla: tamaño, faltantes y duplicados
auditar <- function(datos, nombre, llave) {
  tibble(
    tabla      = nombre,
    filas      = nrow(datos),
    columnas   = ncol(datos),
    celdas_na  = sum(is.na(datos)),               # valores faltantes
    filas_dup  = sum(duplicated(datos)),          # filas idénticas
    llaves_dup = sum(duplicated(datos[[llave]]))  # llave repetida
  )
}

auditoria <- bind_rows(
  auditar(transacciones, "transacciones", "transaccion_id"),
  auditar(clientes,      "clientes",      "cliente_id"),
  auditar(productos,     "productos",     "producto_id"),
  auditar(tiendas,       "tiendas",       "tienda_id"),
  auditar(empleados,     "empleados",     "empleado_id")
)
auditoria
\end{rcode}

\begin{rsalida}
# A tibble: 5 x 6
  tabla         filas columnas celdas_na filas_dup llaves_dup
  <chr>         <int>    <int>     <int>     <int>      <int>
1 transacciones  1000       14         0         0          0
2 clientes        200       10         0         0          0
3 productos        50       13         0         0          0
4 tiendas          10        8         0         0          0
5 empleados       100       12         0         0          0
\end{rsalida}

Buenas noticias en lo estructural: ninguna tabla tiene valores faltantes,
filas duplicadas ni llaves repetidas. Eso no significa que los datos sean
correctos; significa que están \emph{completos}. Sigue con los rangos y la
consistencia interna de la tabla de hechos.

\begin{rcode}
# Rangos de las columnas clave de la tabla de hechos
rango <- function(x) str_c(min(x), " a ", max(x))
transacciones %>%
  summarise(fecha    = rango(fecha),
            hora     = rango(hms::as_hms(hora)),
            cantidad = rango(cantidad),
            precio   = rango(precio_venta),
            total    = rango(total_transaccion)) %>%
  pivot_longer(everything(), names_to = "columna",
               values_to = "rango")

# ¿El total cuadra con cantidad x precio? (tolerancia de 1 centavo)
transacciones %>%
  summarise(totales_incorrectos =
              sum(abs(total_transaccion - cantidad * precio_venta) > 0.01))
\end{rcode}

\begin{rsalida}
# A tibble: 5 x 2
  columna  rango                  
  <chr>    <chr>                  
1 fecha    2023-01-01 a 2023-12-31
2 hora     28800 a 74700          
3 cantidad 1 a 5                  
4 precio   20.14 a 598.81         
5 total    22.65 a 2992.35        
# A tibble: 1 x 1
  totales_incorrectos
                <int>
1                   0
\end{rsalida}

Las fechas cubren el año completo, las horas van de 8:00 a 20:45 (horario de
tienda razonable), no hay cantidades ni precios negativos y los 1,000 totales
cuadran con su detalle. Ahora revisa la \termino{integridad referencial}:
toda llave de la tabla de hechos debe existir en su catálogo. Para eso
\codigo{anti\_join()} es ideal: devuelve las filas de la primera tabla que
\emph{no} encuentran pareja en la segunda (Módulo~\ref{cap:manipulacion}).

\begin{rcode}
# Llaves huérfanas: transacciones cuyo id no existe en el catálogo
tibble(
  llave = c("cliente_id", "producto_id", "tienda_id", "vendedor_id"),
  huerfanas = c(
    nrow(anti_join(transacciones, clientes,  by = "cliente_id")),
    nrow(anti_join(transacciones, productos, by = "producto_id")),
    nrow(anti_join(transacciones, tiendas,   by = "tienda_id")),
    nrow(anti_join(transacciones, empleados,
                   by = c("vendedor_id" = "empleado_id")))
  )
)

# En sentido contrario: elementos del catálogo sin ninguna venta
clientes_sin_compra <- anti_join(clientes, transacciones,
                                 by = "cliente_id")
clientes_sin_compra %>% select(cliente_id, nombre, fecha_registro)
nrow(anti_join(productos, transacciones, by = "producto_id"))
\end{rcode}

\begin{rsalida}
# A tibble: 4 x 2
  llave       huerfanas
  <chr>           <int>
1 cliente_id          0
2 producto_id         0
3 tienda_id           0
4 vendedor_id         0
# A tibble: 2 x 3
  cliente_id nombre            fecha_registro
       <dbl> <chr>             <date>        
1        171 Rosa Jiménez      2022-08-27    
2        199 Ricardo Hernández 2023-09-16    
[1] 0
\end{rsalida}

No hay llaves huérfanas: cada transacción apunta a un cliente, producto,
tienda y vendedor que existen. En sentido contrario, 2 de los 200 clientes
registrados no compraron en 2023 (los usaremos para distinguir ``clientes
registrados'' de ``clientes activos'') y los 50 productos tuvieron al menos
una venta.

La parte más valiosa de la auditoría son las \termino{reglas de negocio}:
condiciones que los datos deberían cumplir según cómo funciona la empresa.
Aquí es donde aparecen los hallazgos interesantes.

\begin{rcode}
# Unimos a cada transacción los datos del producto y del cliente
revision <- transacciones %>%
  left_join(productos, by = "producto_id") %>%
  left_join(select(clientes, cliente_id, fecha_registro),
            by = "cliente_id")

vendedores_ids <- unique(transacciones$vendedor_id)

hallazgos <- tibble(
  regla = c(
    "Productos con margen de catálogo negativo",
    "Transacciones con precio menor al costo",
    "Transacciones a +/-20% del precio de catálogo",
    "Productos inactivos que sí se vendieron",
    "Transacciones de productos inactivos",
    "Transacciones anteriores al registro del cliente",
    "Clientes con compras antes de registrarse",
    "Vendedores que no son del depto. de Ventas"),
  casos = c(
    sum(productos$margen_pct < 0),
    sum(revision$precio_venta < revision$costo),
    sum(abs(revision$precio_venta / revision$precio_catalogo - 1)
        <= 0.20),
    n_distinct(revision$producto_id[!revision$activo]),
    sum(!revision$activo),
    sum(revision$fecha < revision$fecha_registro),
    n_distinct(revision$cliente_id[revision$fecha <
                                     revision$fecha_registro]),
    empleados %>%
      filter(empleado_id %in% vendedores_ids,
             departamento != "Ventas") %>% nrow())
)
hallazgos
\end{rcode}

\begin{rsalida}
# A tibble: 8 x 2
  regla                                            casos
  <chr>                                            <int>
1 Productos con margen de catálogo negativo           10
2 Transacciones con precio menor al costo            252
3 Transacciones a +/-20% del precio de catálogo      149
4 Productos inactivos que sí se vendieron              7
5 Transacciones de productos inactivos               137
6 Transacciones anteriores al registro del cliente   114
7 Clientes con compras antes de registrarse           34
8 Vendedores que no son del depto. de Ventas          21
\end{rsalida}

Cada renglón es una conversación con el área de negocio. Revisemos los más
importantes con detalle.

\begin{rcode}
# Detalle: productos cuyo costo supera al precio de catálogo
productos %>%
  filter(margen_pct < 0) %>%
  select(producto_id, categoria, precio_catalogo, costo, margen_pct,
         activo) %>%
  arrange(margen_pct)

# ¿Qué tan lejos está el precio de venta del de catálogo?
revision %>%
  mutate(razon_precio = precio_venta / precio_catalogo) %>%
  summarise(minimo  = min(razon_precio),
            mediana = median(razon_precio),
            maximo  = max(razon_precio))

# ¿A qué departamento pertenecen los vendedores según Recursos Humanos?
empleados %>%
  filter(empleado_id %in% vendedores_ids) %>%
  count(departamento, sort = TRUE)
\end{rcode}

\begin{rsalida}
# A tibble: 10 x 6
   producto_id categoria   precio_catalogo costo margen_pct activo
         <dbl> <chr>                 <dbl> <dbl>      <dbl> <lgl> 
 1           5 Alimentos              25.5 355.     -1292.  TRUE  
 2          46 Automotriz             64.5 178.      -175.  FALSE 
 3          32 Hogar                 143.  210.       -46.7 TRUE  
 4          26 Electrónica           115.  159.       -38.3 TRUE  
 5          29 Belleza               298.  392.       -31.6 TRUE  
 6           9 Hogar                 250.  328.       -31.2 TRUE  
 7           2 Belleza               313.  399.       -27.7 TRUE  
 8          12 Mascotas              289.  369.       -27.4 TRUE  
 9          40 Electrónica            50.1  58.2      -16.0 TRUE  
10          13 Ropa                  348.  393.       -12.8 FALSE 
# A tibble: 1 x 3
  minimo mediana maximo
   <dbl>   <dbl>  <dbl>
1 0.0288   0.750   23.4
# A tibble: 7 x 2
  departamento            n
  <chr>               <int>
1 Atención al Cliente    10
2 RRHH                    4
3 Ventas                  4
4 Marketing               3
5 Finanzas                2
6 IT                      1
7 Operaciones             1
\end{rsalida}

Esto es lo que encontró la auditoría, traducido a lenguaje de negocio:

\begin{itemize}
  \item \textbf{Márgenes negativos en el catálogo.} 10 de los 50 productos
        cuestan más de lo que dice su precio de lista. El caso extremo es el
        producto 5 (Alimentos): precio de \$25.47 y costo de \$354.63, un
        margen de $-1{,}292$\%. Casi seguro es un error de captura (un costo
        por caja registrado como costo unitario, por ejemplo), pero debe
        confirmarlo Compras.
  \item \textbf{El precio de venta no respeta el catálogo.} Solo
        15\% (149) de las transacciones se vendió a menos de 20\% de
        distancia del precio de lista; hay ventas a 3\% del precio de
        catálogo y otras a 23 veces ese precio. O el catálogo está
        desactualizado, o el punto de venta permite capturar cualquier
        precio. En cualquier caso, \emph{el margen que calculemos con el
        costo del catálogo será una estimación}, y en 252 transacciones
        (una de cada cuatro) el precio cobrado quedó por debajo del costo.
  \item \textbf{Productos inactivos con ventas.} Los 7 productos marcados
        como inactivos suman 137 transacciones en 2023. Lo más probable es
        que la bandera \codigo{activo} describa el estado \emph{actual}
        (productos descontinuados después), no el de todo el año.
  \item \textbf{Compras antes del registro.} 114 transacciones de 34 clientes
        ocurrieron antes de su fecha de registro. La fecha de registro no es
        confiable para medir antigüedad; usaremos la fecha de la primera
        compra.
  \item \textbf{Los ``vendedores'' no son de Ventas.} De los 25
        \codigo{vendedor\_id}, solo 4 corresponden a empleados del
        departamento de Ventas según Recursos Humanos. O la llave
        \codigo{vendedor\_id} no es el \codigo{empleado\_id}, o la tabla de
        empleados está desactualizada. Además, la columna \codigo{sucursal}
        de empleados usa nombres (``Matriz'', ``Sucursal Norte''\ldots) que
        no corresponden a las 10 tiendas, así que no podemos asignar
        empleados a tiendas con esa tabla.
\end{itemize}

\begin{negocio}{La auditoría también es un entregable}
En una consultoría real, la lista de hallazgos de calidad suele ser lo
primero que se presenta, y muchas veces es lo más valorado: la dirección
descubre que su catálogo tiene costos imposibles o que el punto de venta
deja cobrar cualquier precio. Presenta cada hallazgo con su evidencia (el
número de casos), su impacto y un responsable de corregirlo.
\end{negocio}

\subsection{Decisiones de limpieza}
\label{sec:m09-decisiones}

Limpiar no es ``borrar lo raro''. Cada decisión se documenta con su
justificación en una \termino{bitácora de decisiones}
(Tabla~\ref{tab:m09-decisiones}), que acompaña al proyecto y se entrega a la
dirección.

\begin{table}[H]
\centering
\small
\caption{Bitácora de decisiones de limpieza.}
\label{tab:m09-decisiones}
\begin{tabularx}{\textwidth}{@{}>{\raggedright\arraybackslash}p{3.6cm} >{\raggedright\arraybackslash}X >{\raggedright\arraybackslash}X@{}}
\toprule
Problema & Decisión & Justificación \\
\midrule
Dos fuentes de ventas distintas & Hechos = \archivo{transacciones.csv}; \archivo{ventas\_retail.csv} solo para descuentos & Más granular, con hora y método de pago; sumar ambas duplicaría ventas \\
10 productos con margen negativo & Se conservan y se marcan con \codigo{margen\_negativo} & Son ventas reales; el error está en el costo y lo debe corregir Compras \\
Precio de venta distinto al de catálogo & El ingreso se calcula con el precio cobrado; el margen se reporta como ``estimado'' & El precio cobrado es lo que entró a caja; el costo del catálogo es dudoso \\
Ventas por debajo del costo & Se conservan con la bandera \codigo{bajo\_costo} & Pueden ser liquidaciones o errores; se reportan para revisión \\
Productos inactivos con ventas & Se conservan con la bandera \codigo{producto\_inactivo} & La bandera describe el estado actual, no el de 2023 \\
Compras antes del registro & No se usa \codigo{fecha\_registro} para antigüedad; se usa la primera compra & La fecha de registro no es confiable \\
Vendedores fuera del depto. de Ventas & Se conserva \codigo{vendedor\_id}; el análisis de salario se marca como preliminar & La llave no está validada con Recursos Humanos \\
Datos personales de clientes & No se copian email ni teléfono al modelo & Minimizar datos personales: el análisis no los necesita \\
\bottomrule
\end{tabularx}
\end{table}

\begin{consejo}
Observa que ninguna decisión borra filas. Borrar transacciones cambiaría el
total de ventas y abriría la discusión de ``qué número es el bueno'' que la
directora quiere terminar. \emph{Marcar} con banderas (columnas lógicas) te
permite excluir casos en un análisis específico sin alterar la verdad
contable.
\end{consejo}

\subsection{Transformar: modelo en estrella}
\label{sec:m09-estrella}

Organizamos los datos en un \termino{modelo en estrella}: una
\termino{tabla de hechos} con las medidas (cantidad, total, costo, margen) y
las llaves, rodeada de \termino{tablas de dimensiones} que describen el
contexto (quién compró, qué, dónde, cuándo y quién vendió). Es el mismo
diseño que usan Power BI, Tableau y los data warehouses corporativos
(Módulo~\ref{cap:bases-datos}). Empezamos por las dimensiones.

\begin{rcode}
# ---- Dimensión cliente (sin datos personales) ----
dim_cliente <- clientes %>%
  select(cliente_id, nombre, edad, genero, ciudad, segmento_edad = segmento,
         es_premium, fecha_registro) %>%
  mutate(con_compras = cliente_id %in% transacciones$cliente_id)

# ---- Dimensión producto (con banderas de calidad) ----
dim_producto <- productos %>%
  select(producto_id, nombre_producto, categoria, subcategoria, marca,
         proveedor, precio_catalogo, costo, margen_pct, activo,
         stock_actual, stock_minimo, estado_stock) %>%
  mutate(margen_negativo = margen_pct < 0)

# ---- Dimensión tienda (antigüedad a la fecha de corte) ----
dim_tienda <- tiendas %>%
  rename(num_empleados = empleados) %>%
  mutate(anios_operacion = floor(time_length(
           interval(fecha_apertura, fecha_corte), "years")))

# ---- Dimensión vendedor (solo los 25 que aparecen en ventas) ----
dim_vendedor <- empleados %>%
  filter(empleado_id %in% vendedores_ids) %>%
  select(vendedor_id = empleado_id, nombre_vendedor = nombre,
         departamento, puesto, salario_mensual, antiguedad_anos) %>%
  mutate(depto_ventas = departamento == "Ventas")

# ---- Dimensión fecha (calendario completo de 2023) ----
dias_es <- c("lunes", "martes", "miércoles", "jueves", "viernes",
             "sábado", "domingo")
dim_fecha <- tibble(fecha = seq(as.Date("2023-01-01"),
                                as.Date("2023-12-31"), by = "day")) %>%
  mutate(anio       = year(fecha),
         mes        = month(fecha),
         trimestre  = quarter(fecha),
         semana_iso = isoweek(fecha),
         dia_semana = factor(dias_es[wday(fecha, week_start = 1)],
                             levels = dias_es),
         fin_semana = wday(fecha, week_start = 1) >= 6)

map_int(list(dim_cliente, dim_producto, dim_tienda, dim_vendedor,
             dim_fecha), nrow)
\end{rcode}

\begin{rsalida}
[1] 200  50  10  25 365
\end{rsalida}

Construimos el día de la semana con un vector propio (\codigo{dias\_es}) en
lugar de \codigo{wday(label = TRUE)}, porque este último depende del idioma
de tu computadora: en una laptop configurada en inglés obtendrías
\emph{Monday} y tu tablero mezclaría idiomas. Ahora la tabla de hechos:

\begin{rcode}
hechos_ventas <- transacciones %>%
  select(transaccion_id, fecha, hora, cliente_id, producto_id,
         tienda_id, vendedor_id, metodo_pago, cantidad, precio_venta,
         total = total_transaccion) %>%
  # Costo unitario vigente y banderas que vienen de los catálogos
  left_join(select(productos, producto_id, costo_unitario = costo,
                   activo), by = "producto_id") %>%
  left_join(select(clientes, cliente_id, fecha_registro),
            by = "cliente_id") %>%
  mutate(
    hora_dia          = hour(hora),                 # 8, 9, ..., 20
    costo_total       = cantidad * costo_unitario,
    margen_estimado   = total - costo_total,
    bajo_costo        = precio_venta < costo_unitario,
    producto_inactivo = !activo,
    antes_registro    = fecha < fecha_registro
  ) %>%
  select(-activo, -fecha_registro)

# Control: un join mal hecho puede duplicar filas sin avisar
stopifnot(nrow(hechos_ventas) == nrow(transacciones))
stopifnot(sum(hechos_ventas$total) == sum(transacciones$total_transaccion))

hechos_ventas %>%
  summarise(filas = n(), ventas = sum(total),
            margen_estimado = sum(margen_estimado),
            bajo_costo = sum(bajo_costo),
            inactivos = sum(producto_inactivo),
            antes_registro = sum(antes_registro))
\end{rcode}

\begin{rsalida}
# A tibble: 1 x 6
  filas  ventas margen_estimado bajo_costo inactivos antes_registro
  <int>   <dbl>           <dbl>      <int>     <int>          <int>
1  1000 959993.         444637.        252       137            114
\end{rsalida}

Los dos \codigo{stopifnot()} son \termino{controles de conciliación}: si
algún \codigo{left\_join()} duplicara filas (porque una llave estuviera
repetida en el catálogo), el script se detendría en lugar de reportar ventas
infladas. La tabla de hechos conserva las 1,000 transacciones y los
\$959,993 de ventas, y las banderas coinciden con la auditoría. La
Figura~\ref{fig:m09-estrella} muestra el modelo resultante.

\begin{figure}[H]
\centering
\begin{tikzpicture}[
    tabla/.style={draw, rounded corners=2pt, align=left, font=\scriptsize\ttfamily,
                  inner sep=4pt, text width=3.5cm},
    titulo/.style={font=\small\sffamily\bfseries, text=white},
    rel/.style={thick, color=gristexto}]
  \node[tabla, draw=acento, fill=acento!6, text width=4.1cm] (h) {%
    {\sffamily\bfseries\color{acento}hechos\_ventas}\\[2pt]
    transaccion\_id\\ fecha (FK)\\ cliente\_id (FK)\\ producto\_id (FK)\\
    tienda\_id (FK)\\ vendedor\_id (FK)\\ hora, hora\_dia, metodo\_pago\\
    cantidad, precio\_venta, total\\ costo\_total, margen\_estimado\\
    bajo\_costo, producto\_inactivo,\\ antes\_registro};
  \node[tabla, draw=primario, fill=primario!5, above left=0.2cm and 1.3cm of h.north west, anchor=south east] (c) {%
    {\sffamily\bfseries\color{primario}dim\_cliente}\\[2pt]
    cliente\_id (PK)\\ edad, genero, ciudad\\ segmento\_edad, es\_premium\\ con\_compras};
  \node[tabla, draw=primario, fill=primario!5, above right=0.2cm and 1.3cm of h.north east, anchor=south west] (p) {%
    {\sffamily\bfseries\color{primario}dim\_producto}\\[2pt]
    producto\_id (PK)\\ categoria, marca\\ precio\_catalogo, costo\\ activo, stock\_actual\\ margen\_negativo};
  \node[tabla, draw=primario, fill=primario!5, left=1.3cm of h] (f) {%
    {\sffamily\bfseries\color{primario}dim\_fecha}\\[2pt]
    fecha (PK)\\ mes, trimestre\\ semana\_iso\\ dia\_semana, fin\_semana};
  \node[tabla, draw=primario, fill=primario!5, right=1.3cm of h] (t) {%
    {\sffamily\bfseries\color{primario}dim\_tienda}\\[2pt]
    tienda\_id (PK)\\ nombre, ciudad, zona\\ tipo, tamano\_m2\\ num\_empleados};
  \node[tabla, draw=primario, fill=primario!5, below=0.5cm of h] (v) {%
    {\sffamily\bfseries\color{primario}dim\_vendedor}\\[2pt]
    vendedor\_id (PK)\\ departamento, puesto\\ salario\_mensual\\ depto\_ventas};
  \draw[rel] (h) -- (c);
  \draw[rel] (h) -- (p);
  \draw[rel] (h) -- (f);
  \draw[rel] (h) -- (t);
  \draw[rel] (h) -- (v);
\end{tikzpicture}
\caption{Modelo en estrella de TechRetail: una tabla de hechos (naranja) y
cinco dimensiones (azul).}
\label{fig:m09-estrella}
\end{figure}

\subsection{Cargar: guardar el modelo}

El resultado de la Fase~1 se guarda en dos formatos. Los archivos
\termino{RDS} conservan los tipos de R (fechas, horas, factores) y son la
entrada de las fases siguientes y del tablero. La base \termino{SQLite}
permite que otras personas o herramientas (Power BI, Excel, Python) consulten
los mismos datos con SQL.

\begin{rcode}
# ---- RDS: una lista con todas las tablas del modelo ----
modelo <- list(hechos_ventas = hechos_ventas,
               dim_cliente   = dim_cliente,
               dim_producto  = dim_producto,
               dim_tienda    = dim_tienda,
               dim_vendedor  = dim_vendedor,
               dim_fecha     = dim_fecha)
saveRDS(modelo, "resultados/modelo_techretail.rds")

# ---- SQLite: fechas y horas como texto para no perder el formato ----
con <- dbConnect(RSQLite::SQLite(), "datos/techretail.sqlite")
walk2(modelo, names(modelo), function(tabla, nombre) {
  tabla <- tabla %>%
    mutate(across(where(is.Date), as.character),
           across(where(hms::is_hms), as.character),
           across(where(is.factor), as.character))
  dbWriteTable(con, nombre, tabla, overwrite = TRUE)
})
dbListTables(con)

# Comprobación con SQL: ventas por zona directamente en la base
dbGetQuery(con, "
  SELECT t.zona, COUNT(*) AS transacciones,
         ROUND(SUM(h.total), 0) AS ventas
  FROM hechos_ventas h
  JOIN dim_tienda t ON h.tienda_id = t.tienda_id
  GROUP BY t.zona
  ORDER BY ventas DESC")
dbDisconnect(con)
\end{rcode}

\begin{rsalida}
[1] "dim_cliente"   "dim_fecha"     "dim_producto"  "dim_tienda"   
[5] "dim_vendedor"  "hechos_ventas"
    zona transacciones ventas
1  Norte           399 388941
2 Centro           290 275878
3  Oeste           109 105216
4    Sur           108 101538
5   Este            94  88420
\end{rsalida}

SQLite guarda las fechas como números si no las conviertes, por eso las
pasamos a texto antes de escribir (Módulo~\ref{cap:bases-datos}). La consulta
de comprobación es una buena costumbre: demuestra que la base quedó bien
cargada y te da el primer resultado de negocio, la zona Norte (4 sucursales)
concentra la mayor parte de las ventas.

\begin{hazlotu}
Agrega a \codigo{dim\_tienda} una columna \codigo{region} que agrupe las
ciudades en ``Centro'' (CDMX, Puebla, Querétaro) y ``Norte-Occidente''
(Guadalajara, Monterrey, Tijuana), vuelve a guardar el RDS y escribe en SQL
la consulta de ventas por región.
\end{hazlotu}

% ============================================================================
\section{Fase 2: KPIs y análisis exploratorio}
\label{sec:m09-kpis}
% ============================================================================

Con el modelo listo, respondemos las preguntas descriptivas. La regla de esta
fase es simple: \textbf{cada tabla o gráfica termina en una frase de negocio
con números}. Las técnicas de agrupación son las del
Módulo~\ref{cap:manipulacion}; las gráficas siguen los principios del
Módulo~\ref{cap:visualizacion}, y las pruebas estadísticas, los del
Módulo~\ref{cap:eda}.

Primero armamos una \termino{vista analítica}: la tabla de hechos con los
atributos de las dimensiones que vamos a usar. Es el equivalente a una
``tabla plana'' de Excel, pero construida desde el modelo, así que todos los
análisis parten de la misma fuente.

\begin{rcode}
ventas <- hechos_ventas %>%
  left_join(select(dim_producto, producto_id, nombre_producto,
                   categoria), by = "producto_id") %>%
  left_join(select(dim_tienda, tienda_id, nombre_tienda, zona, tipo,
                   tamano_m2), by = "tienda_id") %>%
  left_join(select(dim_cliente, cliente_id, es_premium),
            by = "cliente_id") %>%
  left_join(select(dim_fecha, fecha, mes, dia_semana), by = "fecha")

stopifnot(nrow(ventas) == nrow(hechos_ventas))   # conciliación
\end{rcode}

\subsection{KPIs globales}
\label{sec:m09-kpis-globales}

Un \termino{KPI} (indicador clave de desempeño) es una métrica que la
dirección sigue para saber si va bien. Calcula los globales de 2023 y
preséntalos ya formateados, como aparecerían en la primera diapositiva.

\begin{rcode}
kpis <- ventas %>%
  summarise(ventas_totales  = sum(total),
            transacciones   = n(),
            ticket_promedio = mean(total),
            unidades        = sum(cantidad),
            clientes        = n_distinct(cliente_id),
            margen          = sum(margen_estimado))

tabla_kpis <- tibble(
  kpi = c("Ventas totales", "Transacciones", "Ticket promedio",
          "Unidades vendidas", "Unidades por transacción",
          "Clientes activos", "Base de clientes activa",
          "Venta anual por cliente activo", "Margen bruto estimado",
          "Margen estimado (%)"),
  valor = c(pesos(kpis$ventas_totales),
            comma(kpis$transacciones),
            dollar(kpis$ticket_promedio, accuracy = 0.01),
            comma(kpis$unidades),
            number(kpis$unidades / kpis$transacciones, accuracy = 0.01),
            comma(kpis$clientes),
            percent(kpis$clientes / nrow(dim_cliente), accuracy = 0.1),
            pesos(kpis$ventas_totales / kpis$clientes),
            pesos(kpis$margen),
            percent(kpis$margen / kpis$ventas_totales, accuracy = 0.1))
)
tabla_kpis
\end{rcode}

\begin{rsalida}
# A tibble: 10 x 2
   kpi                            valor   
   <chr>                          <chr>   
 1 Ventas totales                 $959,993
 2 Transacciones                  1,000   
 3 Ticket promedio                $959.99 
 4 Unidades vendidas              2,969   
 5 Unidades por transacción       2.97    
 6 Clientes activos               198     
 7 Base de clientes activa        99.0%   
 8 Venta anual por cliente activo $4,848  
 9 Margen bruto estimado          $444,637
10 Margen estimado (%)            46.3%   
\end{rsalida}

\textbf{Hallazgo:} en 2023 TechRetail vendió \$959,993 en 1,000
transacciones, con un ticket promedio de \$959.99 y casi 3 unidades por
compra. 198 de los 200 clientes registrados compraron al menos una vez
(99\%), y cada cliente activo dejó en promedio \$4,848 en el año. El margen
bruto estimado es de 46.3\%, pero recuerda la auditoría: con el costo del
catálogo en duda, ese número debe presentarse como \emph{estimado}.

La Tabla~\ref{tab:m09-diccionario} es el \termino{diccionario de KPIs}: la
definición exacta de cada indicador. Parece burocracia, pero evita la mitad
de las discusiones de las juntas (``¿ticket promedio por transacción o por
cliente?'').

\begin{table}[H]
\centering
\small
\caption{Diccionario de KPIs del proyecto.}
\label{tab:m09-diccionario}
\begin{tabularx}{\textwidth}{@{}l X l@{}}
\toprule
KPI & Definición & Fuente \\
\midrule
Ventas totales & Suma de \codigo{total} (precio cobrado $\times$ cantidad) & hechos \\
Ticket promedio & Ventas totales / número de transacciones & hechos \\
Clientes activos & Clientes distintos con al menos una compra en el periodo & hechos \\
Base activa & Clientes activos / clientes registrados & hechos + dim\_cliente \\
Margen estimado & Ventas $-$ cantidad $\times$ costo unitario del catálogo & hechos + dim\_producto \\
Crecimiento mensual & Ventas del mes / ventas del mes anterior $-$ 1 & hechos + dim\_fecha \\
Ventas por m\textsuperscript{2} & Ventas de la sucursal / \codigo{tamano\_m2} & hechos + dim\_tienda \\
\bottomrule
\end{tabularx}
\end{table}

\subsection{Evolución mensual}
\label{sec:m09-mensual}

\begin{rcode}
meses_es <- c("ene", "feb", "mar", "abr", "may", "jun",
              "jul", "ago", "sep", "oct", "nov", "dic")

ventas_mes <- ventas %>%
  group_by(mes) %>%
  summarise(ventas = sum(total), transacciones = n(),
            ticket = mean(total), .groups = "drop") %>%
  mutate(crecimiento = ventas / lag(ventas) - 1,          # vs mes previo
         nombre_mes  = factor(meses_es[mes], levels = meses_es))

ventas_mes %>%
  transmute(nombre_mes, ventas = round(ventas), transacciones,
            ticket = round(ticket), crecimiento = percent(crecimiento, 0.1))
\end{rcode}

\begin{rsalida}
# A tibble: 12 x 5
   nombre_mes ventas transacciones ticket crecimiento
   <fct>       <dbl>         <int>  <dbl> <chr>      
 1 ene         82560            98    842 <NA>       
 2 feb         65899            79    834 -20.2%     
 3 mar         90931            91    999 38.0%      
 4 abr         57038            71    803 -37.3%     
 5 may         93910            88   1067 64.6%      
 6 jun         88509            84   1054 -5.8%      
 7 jul         86456            83   1042 -2.3%      
 8 ago         78929            75   1052 -8.7%      
 9 sep         68856            86    801 -12.8%     
10 oct         70004            74    946 1.7%       
11 nov         80507            80   1006 15.0%      
12 dic         96394            91   1059 19.7%      
\end{rsalida}

\begin{rcode}
promedio_mes <- mean(ventas_mes$ventas)
extremos <- ventas_mes %>% filter(ventas %in% range(ventas))

ggplot(ventas_mes, aes(x = nombre_mes, y = ventas, group = 1)) +
  geom_hline(yintercept = promedio_mes, linetype = "dashed",
             color = color_gris, linewidth = 0.6) +
  geom_line(color = color_principal, linewidth = 0.8) +
  geom_point(color = color_principal, size = 2) +
  geom_point(data = extremos, color = color_resalte, size = 3) +
  geom_text(data = extremos, aes(label = pesos(ventas)),
            vjust = -1.1, color = color_resalte, size = 3.5) +
  annotate("text", x = 1, y = promedio_mes, vjust = -0.6, hjust = 0.2,
           label = "Promedio", color = "#52514e", size = 3) +
  scale_y_continuous(labels = dollar, limits = c(0, 110000)) +
  labs(title = "Ventas mensuales de TechRetail, 2023",
       subtitle = "Mucha variación de un mes a otro, sin tendencia clara",
       x = NULL, y = "Ventas") +
  tema_libro()
\end{rcode}

\begin{figure}[H]
\centering
\includegraphics[width=0.9\textwidth]{m09-ventas-mensuales}
\caption{Ventas mensuales con el mejor y el peor mes resaltados.}
\label{fig:m09-ventas-mensuales}
\end{figure}

\textbf{Hallazgo:} diciembre fue el mejor mes (\$96,394) y abril el peor
(\$57,038). El crecimiento mensual oscila entre $-37.3$\% (abril) y
$+64.6$\% (mayo): los saltos son grandes pero no forman una tendencia; la
línea sube y baja alrededor de un promedio de \$80,000 al mes. Para la
dirección, esto significa que \emph{no hay evidencia de crecimiento ni de
caída en 2023} y que comparar un mes contra el anterior genera falsas
alarmas. Es mejor comparar contra el promedio o contra el mismo mes del año
anterior (cuando exista).

\begin{advertencia}
Con un solo año no puedes afirmar que ``diciembre es temporada alta''. Un
diciembre fuerte en un año puede ser azar; la estacionalidad se confirma
cuando el patrón se repite al menos dos o tres años
(Sección~\ref{sec:m09-pronostico}).
\end{advertencia}

\subsection{Categorías: ventas contra margen}
\label{sec:m09-categorias}

\begin{rcode}
ventas_categoria <- ventas %>%
  group_by(categoria) %>%
  summarise(ventas = sum(total), margen = sum(margen_estimado),
            .groups = "drop") %>%
  mutate(participacion = ventas / sum(ventas),
         pct_margen    = margen / ventas) %>%
  arrange(desc(ventas))

ventas_categoria %>%
  mutate(across(c(ventas, margen), round),
         across(c(participacion, pct_margen), ~ percent(.x, 0.1)))
\end{rcode}

\begin{rsalida}
# A tibble: 10 x 5
   categoria   ventas margen participacion pct_margen
   <chr>        <dbl>  <dbl> <chr>         <chr>     
 1 Belleza     182143  56054 19.0%         30.8%     
 2 Electrónica 132672  73067 13.8%         55.1%     
 3 Juguetes    125093  57714 13.0%         46.1%     
 4 Alimentos   121543  73870 12.7%         60.8%     
 5 Automotriz  104085  55538 10.8%         53.4%     
 6 Deportes     86171  48301 9.0%          56.1%     
 7 Hogar        63407  13079 6.6%          20.6%     
 8 Mascotas     57668  35423 6.0%          61.4%     
 9 Libros       57162  25757 6.0%          45.1%     
10 Ropa         30048   5835 3.1%          19.4%     
\end{rsalida}

\begin{rcode}
ventas_categoria %>%
  mutate(grupo = if_else(pct_margen < 0.25, "Margen estimado < 25%",
                         "Margen estimado de 25% o más"),
         etiqueta = if_else(pct_margen < 0.25,
                            str_c("margen ", percent(pct_margen, 1)), "")) %>%
  ggplot(aes(x = ventas, y = fct_reorder(categoria, ventas),
             fill = grupo)) +
  geom_col(width = 0.7) +
  geom_text(aes(label = etiqueta), hjust = -0.1, size = 3.3,
            color = color_resalte) +
  scale_fill_manual(values = c(color_resalte, color_gris)) +
  scale_x_continuous(labels = dollar, expand = expansion(c(0, 0.18))) +
  labs(title = "Ventas por categoría, 2023",
       subtitle = "Belleza lidera en ventas, pero con margen bajo",
       x = "Ventas", y = NULL, fill = NULL) +
  tema_libro()
\end{rcode}

\begin{figure}[H]
\centering
\includegraphics[width=0.9\textwidth]{m09-categorias}
\caption{Ventas por categoría; en naranja, las de margen estimado bajo.}
\label{fig:m09-categorias}
\end{figure}

\textbf{Hallazgo:} Belleza es la categoría que más vende (19\% de las
ventas), pero su margen estimado es de solo 30.8\% por dos productos con
costo mayor al precio de catálogo. Hogar y Ropa tienen los márgenes más bajos
(alrededor de 20\%). En el otro extremo, Mascotas y Alimentos rinden más de
60\% de margen. \textbf{Implicación:} crecer en Belleza sin corregir costos o
precios sería vender más para ganar proporcionalmente menos; la primera
acción es revisar los productos con margen negativo de la auditoría.

\subsection{Sucursales y zonas}
\label{sec:m09-tiendas}

\begin{rcode}
ventas_tienda <- ventas %>%
  group_by(nombre_tienda, zona, tipo, tamano_m2) %>%
  summarise(ventas = sum(total), transacciones = n(),
            ticket = mean(total), .groups = "drop") %>%
  arrange(desc(ventas))

ventas_tienda %>% mutate(across(c(ventas, ticket), round))

# Por zona: total y promedio por sucursal (las zonas tienen
# distinto número de tiendas)
ventas_tienda %>%
  group_by(zona) %>%
  summarise(tiendas = n(), ventas = round(sum(ventas)),
            ventas_por_tienda = round(mean(ventas)), .groups = "drop") %>%
  arrange(desc(ventas_por_tienda))
\end{rcode}

\begin{rsalida}
# A tibble: 10 x 7
   nombre_tienda    zona   tipo  tamano_m2 ventas transacciones ticket
   <chr>            <chr>  <chr>     <dbl>  <dbl>         <int>  <dbl>
 1 Sucursal Outlet  Norte  Outl~       350 114377           106   1079
 2 Sucursal Sur     Norte  Está~       600 105305           109    966
 3 Sucursal Plaza A Oeste  Prin~       700 105216           109    965
 4 Sucursal Express Centro Expr~       300 102791           108    952
 5 Sucursal Premium Sur    Prem~       800 101538           108    940
 6 Sucursal Norte   Norte  Está~       450  89024            99    899
 7 Sucursal Este    Este   Está~       400  88420            94    941
 8 Sucursal Centro  Centro Prin~       500  88414            87   1016
 9 Sucursal Oeste   Centro Está~       550  84672            95    891
10 Sucursal Plaza B Norte  Está~       480  80236            85    944
# A tibble: 5 x 4
  zona   tiendas ventas ventas_por_tienda
  <chr>    <int>  <dbl>             <dbl>
1 Norte        4 388941            388941
2 Centro       3 275878            275878
3 Oeste        1 105216            105216
4 Sur          1 101538            101538
5 Este         1  88420             88420
\end{rsalida}

\begin{rcode}
ventas_tienda %>%
  mutate(tienda  = str_c(nombre_tienda, " (", tamano_m2, " m²)"),
         formato = if_else(tamano_m2 < 400, "Menos de 400 m²",
                           "400 m² o más")) %>%
  ggplot(aes(x = ventas, y = fct_reorder(tienda, ventas),
             fill = formato)) +
  geom_col(width = 0.7) +
  scale_fill_manual(values = c("Menos de 400 m²" = color_resalte,
                               "400 m² o más" = color_gris)) +
  scale_x_continuous(labels = dollar, expand = expansion(c(0, 0.05))) +
  labs(title = "Ventas por sucursal, 2023",
       subtitle = "Las dos tiendas más chicas están entre las 4 que más venden",
       x = "Ventas", y = NULL, fill = "Tamaño") +
  tema_libro()
\end{rcode}

\begin{figure}[H]
\centering
\includegraphics[width=0.9\textwidth]{m09-tiendas}
\caption{Ventas por sucursal; en naranja, los formatos de menos de
400~m\textsuperscript{2}.}
\label{fig:m09-tiendas}
\end{figure}

\textbf{Hallazgo:} la sucursal que más vende es la Outlet de Tijuana
(\$114,377), una de las más pequeñas (350~m\textsuperscript{2}), y la
Express de Querétaro, la más chica de todas, es cuarta. La Premium, con
800~m\textsuperscript{2} y 40 empleados, queda en quinto lugar. Entre la
mejor y la peor sucursal (Plaza B, \$80,236) hay una diferencia de 43\%. Por
zona, la Norte vende más en total (\$388,941) solo porque tiene 4 tiendas;
por sucursal, la zona Oeste (una tienda) es la más alta. \textbf{Implicación:}
el tamaño no explica las ventas; la pregunta correcta es la productividad por
m\textsuperscript{2} y por empleado (Reto~\ref{ej:p-2}), clave si la
dirección evalúa abrir más formatos chicos.

\subsection{Método de pago}
\label{sec:m09-pago}

\begin{rcode}
ventas %>%
  group_by(metodo_pago) %>%
  summarise(transacciones = n(), ventas = sum(total),
            ticket = mean(total), .groups = "drop") %>%
  mutate(pct_ventas = percent(ventas / sum(ventas), 0.1),
         across(c(ventas, ticket), round)) %>%
  arrange(desc(ventas))

# ¿El ticket difiere según el método de pago? (Módulo 4)
kruskal.test(total ~ metodo_pago, data = ventas)$p.value
\end{rcode}

\begin{rsalida}
# A tibble: 4 x 5
  metodo_pago     transacciones ventas ticket pct_ventas
  <chr>                   <int>  <dbl>  <dbl> <chr>     
1 Tarjeta Crédito           350 351171   1003 36.6%     
2 Efectivo                  315 303735    964 31.6%     
3 Tarjeta Débito            233 220855    948 23.0%     
4 Transferencia             102  84232    826 8.8%      
[1] 0.2027074
\end{rsalida}

\textbf{Hallazgo:} la tarjeta de crédito es el método más usado (350
transacciones, 36.6\% de las ventas) y tiene el ticket más alto (\$1,003);
la transferencia es el menos usado (8.8\%) y su ticket es el más bajo
(\$826). Sin embargo, la prueba de Kruskal-Wallis da un valor~$p$ de
P_KRUSKAL: la diferencia de tickets \emph{no es estadísticamente
significativa}. \textbf{Implicación:} no justifica, por sí sola, promociones
para empujar el pago con crédito; sí vale la pena revisar las comisiones
bancarias, porque 61\% de las ventas se paga con tarjeta.

\subsection{Días y horas}
\label{sec:m09-dia-hora}

\begin{rcode}
# Ventas por día de la semana (dia_semana ya es un factor ordenado)
ventas %>%
  group_by(dia_semana) %>%
  summarise(transacciones = n(), ventas = round(sum(total)),
            .groups = "drop")

# Mapa de calor día x hora
calor <- ventas %>%
  count(dia_semana, hora_dia, wt = total, name = "ventas")

ggplot(calor, aes(x = factor(hora_dia), y = fct_rev(dia_semana),
                  fill = ventas)) +
  geom_tile(color = "white", linewidth = 0.6) +
  scale_fill_gradient(low = "#e3eefb", high = "#0f4c92",
                      labels = dollar) +
  labs(title = "Ventas por día de la semana y hora, 2023",
       subtitle = "Sin horas pico claras: las celdas oscuras están dispersas",
       x = "Hora del día", y = NULL, fill = "Ventas") +
  tema_libro() +
  theme(panel.grid.major = element_blank(),
        legend.key.width = unit(1.5, "cm"))
\end{rcode}

\begin{rsalida}
# A tibble: 7 x 3
  dia_semana transacciones ventas
  <fct>              <int>  <dbl>
1 lunes                136 128878
2 martes               129 119364
3 miércoles            155 141329
4 jueves               146 145280
5 viernes              158 161266
6 sábado               130 115172
7 domingo              146 148703
\end{rsalida}

\begin{figure}[H]
\centering
\includegraphics[width=0.9\textwidth]{m09-dia-hora}
\caption{Mapa de calor de ventas por día y hora (escala secuencial de un
solo tono).}
\label{fig:m09-dia-hora}
\end{figure}

\textbf{Hallazgo:} el viernes es el día con más ventas (\$161,266) y,
contra lo que dicta la intuición del retail, el sábado es el de menos
(\$115,172). El mapa de calor no muestra una franja horaria dominante: las
celdas más oscuras aparecen dispersas. Ojo con la escala: 1,000
transacciones repartidas en 91 celdas (7 días $\times$ 13 horas) son unas 11
por celda, así que una sola venta grande pinta una celda de oscuro.
\textbf{Implicación:} antes de mover turnos de personal con base en este
mapa, conviene acumular más meses de datos; con estos números no hay una
hora pico que justifique reforzar el personal.

\subsection{Pareto de productos}
\label{sec:m09-pareto}

El \termino{principio de Pareto} (80/20) dice que en muchos negocios cerca
del 20\% de los productos genera el 80\% de las ventas. ¿Se cumple en
TechRetail?

\begin{rcode}
pareto <- ventas %>%
  group_by(producto_id, nombre_producto, categoria) %>%
  summarise(ventas = sum(total), .groups = "drop") %>%
  arrange(desc(ventas)) %>%
  mutate(rango         = row_number(),
         pct_productos = rango / n(),
         pct_acumulado = cumsum(ventas) / sum(ventas))

pareto %>%
  summarise(
    ventas_top20pct  = percent(max(pct_acumulado[pct_productos <= 0.2]),
                               0.1),
    productos_al_80  = min(rango[pct_acumulado >= 0.8]),
    mayor_vs_menor   = round(max(ventas) / min(ventas), 1))
\end{rcode}

\begin{rsalida}
# A tibble: 1 x 3
  ventas_top20pct productos_al_80 mayor_vs_menor
  <chr>                     <int>          <dbl>
1 26.4%                        37            3.3
\end{rsalida}

\begin{rcode}
ggplot(pareto, aes(x = pct_productos, y = pct_acumulado)) +
  geom_abline(slope = 1, intercept = 0, linetype = "dashed",
              color = color_gris, linewidth = 0.6) +
  geom_hline(yintercept = 0.8, linetype = "dotted", color = color_gris) +
  geom_vline(xintercept = 0.2, linetype = "dotted", color = color_gris) +
  geom_line(color = color_principal, linewidth = 0.8) +
  annotate("point", x = 0.2, y = 0.264, color = color_resalte, size = 3) +
  annotate("text", x = 0.22, y = 0.23, hjust = 0, size = 3.4,
           color = color_resalte,
           label = "El 20% de los productos genera 26% de las ventas") +
  annotate("text", x = 0.62, y = 0.52, hjust = 0, size = 3,
           color = "#52514e", label = "Reparto perfectamente igual") +
  scale_x_continuous(labels = percent) +
  scale_y_continuous(labels = percent) +
  labs(title = "Curva de Pareto de productos, 2023",
       subtitle = "Las ventas están repartidas: no hay un grupo de productos estrella",
       x = "Productos (ordenados de mayor a menor venta)",
       y = "Ventas acumuladas") +
  tema_libro()
\end{rcode}

\begin{figure}[H]
\centering
\includegraphics[width=0.9\textwidth]{m09-pareto}
\caption{Curva de Pareto: la curva azul está cerca de la diagonal de reparto
igual.}
\label{fig:m09-pareto}
\end{figure}

\textbf{Hallazgo:} los 10 productos más vendidos (20\% del catálogo) generan
solo 26.4\% de las ventas, y se necesitan 37 productos (74\% del catálogo)
para llegar al 80\%. El producto que más vende factura 3.3 veces lo que el
que menos vende. \textbf{Implicación:} TechRetail \emph{no} depende de unos
cuantos productos estrella; una estrategia de ``concentrar el inventario en
el top 10'' dejaría fuera casi tres cuartas partes de las ventas. Es un
resultado ``negativo'' igual de valioso que uno positivo: descarta una
estrategia.

\subsection{Desempeño de vendedores}
\label{sec:m09-vendedores}

\begin{rcode}
ventas_vendedor <- ventas %>%
  group_by(vendedor_id) %>%
  summarise(ventas = sum(total), transacciones = n(),
            .groups = "drop") %>%
  left_join(dim_vendedor, by = "vendedor_id") %>%
  arrange(desc(ventas))

ventas_vendedor %>%
  select(vendedor_id, departamento, puesto, ventas, transacciones,
         salario_mensual) %>%
  mutate(ventas = round(ventas)) %>%
  slice(c(1:3, 23:25))           # los 3 mejores y los 3 últimos

# ¿Venden más quienes ganan más? (correlación, Módulo 4)
prueba_salario <- cor.test(ventas_vendedor$ventas,
                           ventas_vendedor$salario_mensual)
c(r = round(unname(prueba_salario$estimate), 3),
  p_valor = round(prueba_salario$p.value, 3))
\end{rcode}

\begin{rsalida}
# A tibble: 6 x 6
  vendedor_id departamento puesto ventas transacciones salario_mensual
        <dbl> <chr>        <chr>   <dbl>         <int>           <dbl>
1           7 RRHH         Coord~  50470            45          14488.
2           2 Atención al~ Asist~  50168            39          15293.
3          11 Atención al~ Coord~  50017            48          30743.
4          14 Ventas       Geren~  26430            33          38782.
5           1 Marketing    Anali~  24281            25          17291.
6          10 IT           Espec~  21678            24          11205.
      r p_valor 
 -0.138   0.510 
\end{rsalida}

\begin{rcode}
ventas_vendedor %>%
  mutate(grupo = if_else(depto_ventas, "Depto. de Ventas",
                         "Otro departamento")) %>%
  ggplot(aes(x = salario_mensual, y = ventas, color = grupo)) +
  geom_point(size = 2.8) +
  scale_color_manual(values = c("Depto. de Ventas" = color_resalte,
                                "Otro departamento" = color_gris)) +
  scale_x_continuous(labels = dollar) +
  scale_y_continuous(labels = dollar) +
  labs(title = "Ventas anuales contra salario mensual por vendedor",
       subtitle = "Sin relación (r = -0.14, p = 0.51); dato de RR. HH. por validar",
       x = "Salario mensual", y = "Ventas 2023", color = NULL) +
  tema_libro()
\end{rcode}

\begin{figure}[H]
\centering
\includegraphics[width=0.9\textwidth]{m09-vendedores}
\caption{Ventas contra salario de los 25 vendedores.}
\label{fig:m09-vendedores}
\end{figure}

\textbf{Hallazgo:} el mejor vendedor (id 7) vendió \$50,470 y el último
(id 10) \$21,678, 2.3 veces menos. La correlación entre ventas y salario es
prácticamente nula ($r = -0.14$, $p = 0.51$): ganar más no se asocia con
vender más. \textbf{Implicación y cautela:} esto sugiere revisar el esquema
de compensación (por ejemplo, comisiones sobre venta), pero recuerda que la
auditoría mostró que 21 de los 25 ``vendedores'' figuran en otros
departamentos. Antes de presentar este hallazgo, valida la llave con Recursos
Humanos; si la llave está mal, la gráfica compara ventas de una persona con
el salario de otra.

\subsection{Clientes premium contra no premium}
\label{sec:m09-premium}

El programa premium tiene un costo (beneficios, descuentos). Si funciona,
los clientes premium deberían gastar más. Es una pregunta de
\termino{prueba de hipótesis} (Módulo~\ref{cap:eda}).

\begin{rcode}
gasto_cliente <- ventas %>%
  group_by(cliente_id, es_premium) %>%
  summarise(gasto = sum(total), compras = n(), .groups = "drop") %>%
  mutate(grupo = if_else(es_premium, "Premium", "No premium"))

gasto_cliente %>%
  group_by(grupo) %>%
  summarise(clientes = n(), gasto_promedio = round(mean(gasto)),
            gasto_mediano = round(median(gasto)),
            compras_promedio = round(mean(compras), 2), .groups = "drop")

# Prueba t de Welch y, como respaldo, Wilcoxon (no paramétrica)
t.test(gasto ~ grupo, data = gasto_cliente)
wilcox.test(gasto ~ grupo, data = gasto_cliente)$p.value
\end{rcode}

\begin{rsalida}
# A tibble: 2 x 5
  grupo      clientes gasto_promedio gasto_mediano compras_promedio
  <chr>         <int>          <dbl>         <dbl>            <dbl>
1 No premium      126           4929          4525             5.18
2 Premium          72           4708          4517             4.82

	Welch Two Sample t-test

data:  gasto by grupo
t = 0.56874, df = 160.59, p-value = 0.5703
alternative hypothesis: true difference in means between group No premium and group Premium is not equal to 0
95 percent confidence interval:
 -544.5031  984.9854
sample estimates:
mean in group No premium    mean in group Premium 
                4928.536                 4708.294 

[1] 0.8882566
\end{rsalida}

\begin{rcode}
ggplot(gasto_cliente, aes(x = grupo, y = gasto)) +
  geom_boxplot(fill = color_principal, alpha = 0.25, width = 0.5,
               color = color_principal, outlier.shape = NA) +
  geom_jitter(width = 0.15, alpha = 0.5, size = 1.5,
              color = color_principal) +
  stat_summary(fun = mean, geom = "point", shape = 18, size = 4,
               color = color_resalte) +
  scale_y_continuous(labels = dollar) +
  labs(title = "Gasto anual por cliente: premium contra no premium",
       subtitle = "Distribuciones casi idénticas (rombo naranja = promedio)",
       x = NULL, y = "Gasto en 2023") +
  tema_libro()
\end{rcode}

\begin{figure}[H]
\centering
\includegraphics[width=0.9\textwidth]{m09-premium}
\caption{Gasto anual por cliente según el programa premium.}
\label{fig:m09-premium}
\end{figure}

\textbf{Hallazgo:} los clientes premium gastaron en promedio GASTO_PREM
contra GASTO_NOPREM de los no premium, y compraron con la misma frecuencia.
La prueba $t$ da $p = $ P_T y la de Wilcoxon $p = $ P_W: \emph{no hay
diferencia significativa}. \textbf{Implicación:} el programa premium, tal
como está, no se traduce en más gasto. Es una recomendación directa para la
dirección: rediseñar los beneficios (atarlos a metas de compra) o medir si
el programa aporta otra cosa (por ejemplo, menor abandono) antes de seguir
financiándolo.

\begin{hazlotu}
Repite la comparación premium contra no premium, pero con el
\emph{ticket promedio} por transacción en lugar del gasto anual. ¿Cambia la
conclusión? Después compara el gasto por \codigo{segmento\_edad} con
\codigo{kruskal.test()}.
\end{hazlotu}

%%FIN_PARTE_2
